\documentclass[fontsize=15pt]{jlreq}

%プリアンブル
\usepackage{amsmath}
\usepackage{mathtools}
\pagestyle{empty}
\usepackage[top=15mm, bottom=20mm, left=20mm, right=20mm]{geometry}

%本文
\begin{document}
\begin{center}
{\LARGE {振り返り}} \\
-電池- %単元ごとに変えるか消す
\end{center}
\vspace{1em}
\begin{flushright}
名前：\underline{\hspace{5cm}}
\end{flushright}

% 問題部分
\begin{enumerate}
    \item 次のイオンをイオン式で答えろ．
    \begin{itemize}
        \item 塩化物イオン\hspace{1.4em}( \hspace{2cm} )
        \item 銅イオン\hspace{3.4em}( \hspace{2cm} )
        \item 水素イオン\hspace{2.4em}( \hspace{2cm} )
        \item 水酸化物イオン\hspace{0.4em}( \hspace{2cm} )
    \end{itemize}
    \vspace{2em} 

    \item 化学変化によって化学エネルギーを電気エネルギーに変える装置をなんというか．\\
    ( \hspace{3cm} )
    \vspace{2em} 

    \item 水に溶かすと電気が流れる物質をなんというか．\\
    ( \hspace{3cm} )
    \vspace{2em} 

    \item $\mathrm{Zn}$，$\mathrm{Mg}$，$\mathrm{Cu}$ をイオンになりやすい順に並べろ．\\
    ( \hspace{2cm}$\rightarrow$ \hspace{2cm} $\rightarrow$\hspace{2cm} )
    \vspace{2em} 

    \item $\mathrm{HCl}$水溶液に銅板とマグネシウム板をおいたとき，陽極はどっちになるか．\\
    ( \hspace{3cm} ) 
    \vspace{2em} 

    \item 次のうち，
    電流が流れる物質はどれか．
    全て選んで記号で答えよ．\\
    \begin{tabular}{ll}
      A：鉄板 & B：塩化ナトリウム \\
      C：塩化ナトリウム水溶液 & D：エタノール \\
    \end{tabular} \\
    (\hspace{4cm})

\end{enumerate}

\end{document}
